\scanpage{14}
\mspage{\handwritten{13}}

$g$ manifestement se prolonge en $\bar{g}:\bar{T}\longrightarrow X\mathbin{x}Y$. Cela nous ramène à \unsure{prouver} que \inserted{le composé} $\handwritten{\unsure{F'|F^i|F^o}}\bar T'\xrightarrow{\handwritten{\alpha_i}} \handwritten{\unsure{F'|F^i|F^o}}Z'_i\xrightarrow{\handwritten{\alpha_i}}\handwritten{F}\bar T'$ ne dépend pas de $i=0,1$. Ecrivons $X$ au lieu de $\bar T'$, \unsure{et} considérons le diagramme cartésien d'inclusions

% The square and its arrowheads are penned; the graph and alpha labels are handwritten.
\diagram{\begin{tikzcd}[ampersand replacement=\&]
X_s\arrow[r,"\handwritten{\alpha_s}"]\arrow[d,"\handwritten{\alpha_s}"'] \& X\arrow[d,"\handwritten{\Gamma\unsure{\triangleq|=|\underline{=}}(p,\mathrm{id}_X)}"]\\
X\arrow[r,"i_s"'] \& S^1\mathbin{x}X\mathpunct{,}
\end{tikzcd}}

on a alors $\handwritten{\alpha}{}_{s}^{\handwritten{*}}\handwritten{\alpha}{}_{s\handwritten{*}}=\Gamma^{\handwritten{*}}i_{s\handwritten{*}}$, qui ne dépend pas de $s$ puisque $i_s$ n'en dépend pas à homotopie près. - Le reste est sans doute l'AQT.

Sur la strcture des $B_{\bullet}(X)$. Ce sont tous des algèbres graduées sur $B_{\bullet}(\mathrm{pt})=\handwritten{\Lambda}$, qui est une algèbre graduée à degrés positifs, réduite à $\underline{F}_2$ en degré\deleted{\unsure{x|s}} $0$. C'est d'ailleurs une algèbre de polynômes connue\inserted{\handwritten{(à une infinité de générateurs, $X_i$ \unsure{$i$ entier|en tout|$i$ entier $\geq 1$} \unsure{$\ne$} \unsure{de la forme} $2$\illegible{1})}}.
\margin{\handwritten{\unsure{?|??}}}

On a un homomorphisme \inserted{d'annaàuxpgradués} évident $B_{\bullet}(X)\longrightarrow H^{\unsure{\mathbb B|B}}(X)$, nul en degrés $\handwritten{>}0$, qui à tout $Z$ sur $X$ associe l'image de $1\handwritten{\in} H^o(Z)$ par l'homomorphisme de \unsure{Gysin} associé à $Z\longrightarrow X$. Cet homomorphisme par Thom est surjectif (toute classe de cohomologie provient d'une image de variété compacte....); \deleted{\unsure{noter}}\unsure{.}
\margin{\handwritten{\uline{NB $H^*(X)$} $(=H^*(X,\mathbb F_2))$}}% slanted
\margin{\handwritten{en fait $H(X)\simeq B_{\bullet}(X)/\Lambda^+B_{\bullet}(X)$}}% slanted

On ppouve même qu'il existe un relèvement fonctoriel en $X$ $H^{\bullet}(X)\longrightarrow B_{\bullet}(X)$, \handwritten{(}\inserted{\handwritten{il est douteux qu'il soit}} compatible avec multiplication\handwritten{)} \deleted{\unsure{kin|kix}}tel que l'on ait:
\margin{\handwritten{?}}

\deleted{a)} L'homomorphisme correspondant $H^{\bullet}(X)\otimes_{\underline F_2}\handwritten{\Lambda\simeq} B_{\bullet}(X)$ est un isomorphisme (ce sera vrai pour tout relèvement individuellement ...), \deleted{et}

b\handwritten{)} \handwritten{[}\inserted{\handwritten{Quillen ignore si}} \deleted{\unsure{Le|Le ce}} relèvement est compatible aussi avec les hom. de Gysin\unsure{.}
\margin{\handwritten{?}}

Ainsi, la catégorie $\underline B$ admet un foncteur pleinement fidèle dans la catégorie des modules gradués libres de type fini sur $\handwritten{\Lambda}$, les modules de type fini qu'on obtient ainsi ayant des générateurs à degrés $\handwritten{<}0$ \unsure{et} satisfaisant (pour $X$ connexe i.e. quand il y a un seul générateur de degré zéro) la dualité de Poincaré. Quand on prend la catégorie \uline{\unsure{quasi} abélienne} $\inserted{\handwritten{\widetilde B}}$ associée à $\underline B$, on trouve donc tous les modules gradués libres de type fini à générateurs \deleted{négatifs} en degrés $\handwritten{\leq}0$ seulement. C'est un
\margin{\handwritten{i.e. l'enveloppe ``Karoubienne''}}% slanted

% archive sheet number: 13
